% GENERATED FILE. Edit canonical lessons, then run npm run generate:books.
% canonical-source-hash: fnv1a64:cd99512a7c7b10d1
% canonical-lessons: PT-C448-tabacaria, PT-C448-isqueiro, PT-C448-fosforo, PT-C448-raspadinha, PT-C448-folhear, PT-R448-revisao-441-448

\chapter{Jornais e tabaco}
\label{ch:pt-jornais-e-tabaco}
\hlchaptermodality{448}

\begin{chapteropening}
\textbf{By the end of this chapter:} \emph{I can use small everyday shops: the tobacconist, a lighter, a match, a scratch card and leafing through a magazine.}

Complete the last lesson of chapter 448: I can use small everyday shops: the tobacconist, a lighter, a match, a scratch card and leafing through a magazine.
\end{chapteropening}

\section[a tabacaria]{a tabacaria — the tobacconist, the newsagent}

\label{lesson:PT-C448-tabacaria}



\begin{quote}
Before the new one: say the Portuguese for to potter about, then the Portuguese for the viewpoint.
\end{quote}

\subsection*{You'll want to know: a tabacaria}
\textbf{a tabacaria} — \textquotedblleft{}the tobacconist, the newsagent\textquotedblright{}.

\begin{quote}
Compro o jornal na tabacaria da esquina.
\end{quote}

\begin{quote}
\emph{I buy the paper at the newsagent on the corner.}
\end{quote}

It sells newspapers, magazines and lottery tickets too.

\begin{cousinweb}
\emph{tabacaria} is built on \emph{tabaco} with \emph{-aria}, a shop (EtymDB); \emph{tabaco} was borrowed from Spanish \emph{tabaco}.
\end{cousinweb}

\begin{grammarlens}[title={What you've built}]
One more word for telling, explaining and describing at B1.
\end{grammarlens}

\subsection*{Your turn}
\begin{itemize}
\raggedright
  \item \emph{Say it:} \emph{a tabacaria}
  \item \emph{Say it:} \emph{a tabacaria}, once more
  \item \emph{Say it:} a sentence of your own with \emph{a tabacaria}
  \item \emph{Recall:} say the Portuguese for to potter about, then the Portuguese for the viewpoint, then say \emph{a tabacaria} again
\end{itemize}

\begin{tcolorbox}[breakable,colback=teal!4,colframe=teal!35!black,title={Before you move on}]
What does a tabacaria mean? (\textquotedblleft{}The tobacconist, the newsagent\textquotedblright{}.) Say it once more.
\end{tcolorbox}
\section[o isqueiro]{o isqueiro — the lighter (for cigarettes)}

\label{lesson:PT-C448-isqueiro}



\begin{quote}
Before the new one: say the Portuguese for the viewpoint, then the Portuguese for the tobacconist.
\end{quote}

\subsection*{You'll want to know: o isqueiro}
\textbf{o isqueiro} — \textquotedblleft{}the lighter (for cigarettes)\textquotedblright{}.

\begin{quote}
Ele perdeu o isqueiro outra vez.
\end{quote}

\begin{quote}
\emph{He lost his lighter again.}
\end{quote}

Older lighters burned a piece of tinder.

\begin{cousinweb}
\emph{isqueiro} is built on \emph{isca} 'bait, tinder' with \emph{-eiro} (EtymDB); \emph{isca} is inherited from Latin \emph{ēsca} 'food, bait'.
\end{cousinweb}

\begin{grammarlens}[title={What you've built}]
One more word for telling, explaining and describing at B1.
\end{grammarlens}

\subsection*{Your turn}
\begin{itemize}
\raggedright
  \item \emph{Say it:} \emph{o isqueiro}
  \item \emph{Say it:} \emph{o isqueiro}, once more
  \item \emph{Say it:} a sentence of your own with \emph{o isqueiro}
  \item \emph{Recall:} say the Portuguese for the viewpoint, then the Portuguese for the tobacconist, then say \emph{o isqueiro} again
\end{itemize}

\begin{tcolorbox}[breakable,colback=teal!4,colframe=teal!35!black,title={Before you move on}]
What does o isqueiro mean? (\textquotedblleft{}The lighter (for cigarettes)\textquotedblright{}.) Say it once more.
\end{tcolorbox}
\section[o fósforo]{o fósforo — the match; phosphorus}

\label{lesson:PT-C448-fosforo}



\begin{quote}
Before the new one: say the Portuguese for the tobacconist, then the Portuguese for the lighter.
\end{quote}

\subsection*{You'll want to know: o fósforo}
\textbf{o fósforo} — \textquotedblleft{}the match; phosphorus\textquotedblright{}.

\begin{quote}
Acendi a vela com um fósforo.
\end{quote}

\begin{quote}
\emph{I lit the candle with a match.}
\end{quote}

\emph{Uma caixa de fósforos} — a box of matches.

\begin{cousinweb}
\emph{fósforo} is a learned borrowing, through Latin, from Greek \emph{phōsphóros} 'light-bearing', from \emph{phôs} 'light' and \emph{phérein} 'to carry' (Lewis \& Short).
\end{cousinweb}

\begin{grammarlens}[title={What you've built}]
One more word for telling, explaining and describing at B1.
\end{grammarlens}

\subsection*{Your turn}
\begin{itemize}
\raggedright
  \item \emph{Say it:} \emph{o fósforo}
  \item \emph{Say it:} \emph{o fósforo}, once more
  \item \emph{Say it:} a sentence of your own with \emph{o fósforo}
  \item \emph{Recall:} say the Portuguese for the tobacconist, then the Portuguese for the lighter, then say \emph{o fósforo} again
\end{itemize}

\begin{tcolorbox}[breakable,colback=teal!4,colframe=teal!35!black,title={Before you move on}]
What does o fósforo mean? (\textquotedblleft{}The match; phosphorus\textquotedblright{}.) Say it once more.
\end{tcolorbox}
\section[a raspadinha]{a raspadinha — the scratch card}

\label{lesson:PT-C448-raspadinha}



\begin{quote}
Before the new one: say the Portuguese for the lighter, then the Portuguese for the match.
\end{quote}

\subsection*{You'll want to know: a raspadinha}
\textbf{a raspadinha} — \textquotedblleft{}the scratch card\textquotedblright{}.

\begin{quote}
A minha tia compra uma raspadinha todas as semanas.
\end{quote}

\begin{quote}
\emph{My aunt buys a scratch card every week.}
\end{quote}

Built on \emph{raspar}, to scratch off.

\begin{cousinweb}
\emph{raspadinha} is a small form built on \emph{raspar} 'to scrape', which is inherited from Vulgar Latin \emph{raspāre} (EtymDB).
\end{cousinweb}

\begin{grammarlens}[title={What you've built}]
One more word for telling, explaining and describing at B1.
\end{grammarlens}

\subsection*{Your turn}
\begin{itemize}
\raggedright
  \item \emph{Say it:} \emph{a raspadinha}
  \item \emph{Say it:} \emph{a raspadinha}, once more
  \item \emph{Say it:} a sentence of your own with \emph{a raspadinha}
  \item \emph{Recall:} say the Portuguese for the lighter, then the Portuguese for the match, then say \emph{a raspadinha} again
\end{itemize}

\begin{tcolorbox}[breakable,colback=teal!4,colframe=teal!35!black,title={Before you move on}]
What does a raspadinha mean? (\textquotedblleft{}The scratch card\textquotedblright{}.) Say it once more.
\end{tcolorbox}
\section[folhear]{folhear — to leaf through, to flick through}

\label{lesson:PT-C448-folhear}



\begin{quote}
Before the new one: say the Portuguese for the match, then the Portuguese for the scratch card.
\end{quote}

\subsection*{You'll want to know: folhear}
\textbf{folhear} — \textquotedblleft{}to leaf through, to flick through\textquotedblright{}.

\begin{quote}
Folheei uma revista enquanto esperava.
\end{quote}

\begin{quote}
\emph{I leafed through a magazine while I waited.}
\end{quote}

Built on \emph{a folha}, the page.

\begin{cousinweb}
\emph{folhear} is built on \emph{folha} 'leaf, page', which is inherited from Latin \emph{folia}, the plural of \emph{folium} 'leaf' (EtymDB).
\end{cousinweb}

\begin{grammarlens}[title={What you've built}]
One more word for telling, explaining and describing at B1.
\end{grammarlens}

\subsection*{Your turn}
\begin{itemize}
\raggedright
  \item \emph{Say it:} \emph{folhear}
  \item \emph{Say it:} \emph{folhear}, once more
  \item \emph{Say it:} a sentence of your own with \emph{folhear}
  \item \emph{Recall:} say the Portuguese for the match, then the Portuguese for the scratch card, then say \emph{folhear} again
\end{itemize}

\begin{tcolorbox}[breakable,colback=teal!4,colframe=teal!35!black,title={Before you move on}]
What does folhear mean? (\textquotedblleft{}To leaf through, to flick through\textquotedblright{}.) Say it once more.
\end{tcolorbox}
\section[(dialogue)]{Revisão dos capítulos 441–448}

\label{lesson:PT-R448-revisao-441-448}



\begin{quote}
Say the Portuguese for the bell tower and the cathedral.
\end{quote}

\subsection*{Your turn}
Say these aloud:

\begin{itemize}
\raggedright
  \item the bell tower, the cathedral, the arcades, the public fountain, the pillory
  \item the city wall, the fortress, the monastery, the cloister, to wander
  \item the police station, the barracks, the registry office, to report, to apply for
  \item the post, the recycling bins, the gutter, the bollard, to obstruct
  \item the scaffolding, the crane, the building site, to collapse, to knock down
  \item the open-air party, the pilgrimage festival, the procession, the patron saint, to parade
  \item the seafront road, to roam, to peek, to potter about, the viewpoint
  \item the tobacconist, the lighter, the match, the scratch card, to leaf through
\end{itemize}

\begin{tcolorbox}[breakable,colback=teal!4,colframe=teal!35!black,title={Before you move on}]
The bell tower? (\textbf{o campanário}.) The cathedral? (\textbf{a catedral}.) The arcades? (\textbf{as arcadas}.) The public fountain? (\textbf{o chafariz}.) The pillory? (\textbf{o pelourinho}.) The city wall? (\textbf{a muralha}.) The fortress? (\textbf{a fortaleza}.) The monastery? (\textbf{o mosteiro}.) The cloister? (\textbf{o claustro}.) To wander? (\textbf{deambular}.) The police station? (\textbf{a esquadra}.) The barracks? (\textbf{o quartel}.) The registry office? (\textbf{a conservatória}.) To report? (\textbf{denunciar}.) To apply for? (\textbf{requerer}.) The post? (\textbf{o poste}.) The recycling bins? (\textbf{o ecoponto}.) The gutter? (\textbf{a sarjeta}.) The bollard? (\textbf{o pilarete}.) To obstruct? (\textbf{obstruir}.) The scaffolding? (\textbf{o andaime}.) The crane? (\textbf{a grua}.) The building site? (\textbf{o estaleiro}.) To collapse? (\textbf{desabar}.) To knock down? (\textbf{demolir}.) The open-air party? (\textbf{o arraial}.) The pilgrimage festival? (\textbf{a romaria}.) The procession? (\textbf{a procissão}.) The patron saint? (\textbf{o padroeiro}.) To parade? (\textbf{desfilar}.) The seafront road? (\textbf{a marginal}.) To roam? (\textbf{vaguear}.) To peek? (\textbf{espreitar}.) To potter about? (\textbf{cirandar}.) The viewpoint? (\textbf{o miradouro}.) The tobacconist? (\textbf{a tabacaria}.) The lighter? (\textbf{o isqueiro}.) The match? (\textbf{o fósforo}.) The scratch card? (\textbf{a raspadinha}.) To leaf through? (\textbf{folhear}.)
\end{tcolorbox}
